\section{A202253 and A202254: bounded zero-sum arrays}

For $r\in\{3,4\}$ and $n\ge0$, let $A_r(n)$ denote the number of
ordered integer arrays $(a_1,\ldots,a_r)$ such that
\[
 -n\le a_i\le n,\qquad \sum_{i=1}^r a_i=0,\qquad
 -n\le a_i-a_{i+1}\le n\quad(1\le i<r).
\]
Thus adjacency is along the linear array, without a wraparound edge.
The source sequences A202253 and A202254 begin at $n=1$
\cite{a202253-source,a202254-source}; at $n=0$ their natural
combinatorial extensions both equal one.

\begin{theorem}\label{thm:bounded-zero-sum}
For every integer $n\ge0$,
\[
 A_3(n)=1+2\left\lfloor\frac{2n^2+2n}{3}\right\rfloor,
\]
and
\[
 A_4(n)=
 \begin{cases}
 (22n^3+33n^2+26n+12)/12,&n\text{ even},\\
 (22n^3+33n^2+26n+3)/12,&n\text{ odd}.
 \end{cases}
\]
These prove the recorded formula conjectures of Kimberling and Barker
for the literal array definitions.
\end{theorem}

\begin{proof}
First consider three entries. Substitute $z=-x-y$ in an array
$(x,y,z)$. For $y\ge0$, the entry bounds and the two adjacent-difference
bounds are equivalent to
\[
 y-n\le x\le n-2y.
\]
The interval is nonempty exactly when $3y\le2n$. Put
$m=\lfloor2n/3\rfloor$. The slice at $y=0$ has $2n+1$ points, and
negation bijects every positive-middle slice with its negative-middle
counterpart. Therefore
\begin{equation}\label{eq:three-slices}
 A_3(n)=2n+1+2\sum_{y=1}^{m}(2n-3y+1)
       =2n+1+m(4n-3m-1).
\end{equation}
For $n=3q,3q+1,3q+2$, this gives respectively
\[
 12q^2+4q+1,\qquad12q^2+12q+3,\qquad12q^2+20q+9.
\]
These are exactly the three residue-class evaluations of the asserted
floor expression.

For four entries, write an array as $(a,b,c,d)$ and put $s=b+c$.
Negation pairs the slices at $s$ and $-s$, so first suppose $s\ge0$.
The two outer differences satisfy
$c-d=(a-b)+2s$; since both lie in $[-n,n]$, we have $0\le s\le n$.
Set
\[
 m=n-s,\qquad h=\lfloor m/2\rfloor,\qquad
 \epsilon=m\bmod2,\qquad v=a+c.
\]
Then
\[
 c=s-b,\qquad a=v-s+b,\qquad d=-v-b.
\]
Substituting these expressions into all the entry and difference
bounds gives the following equivalent conditions:
\begin{equation}\label{eq:four-rectangle}
 -h\le b\le n-h-\epsilon,\qquad
 -m\le v\le m,\qquad -m\le b+v\le n.
\end{equation}
The first two intervals contain
$(n+1-\epsilon)(2m+1)$ integer points. To count the upper forbidden
corner, write
\[
 b=n-h-\epsilon-i,\qquad v=m-j.
\]
Since $m=2h+\epsilon$, the inequality $b+v>n$ becomes $i+j<h$.
For $i,j\ge0$ this is a triangle with $h(h+1)/2$ points. All those
points lie in the rectangle: $m\le n$ ensures the horizontal bound,
and the vertical bound follows from $i+j<h\le m$. Conversely,
every rectangle point above the diagonal gives such a pair $(i,j)$.
The map
\[
 (b,v)\longmapsto(n-m-b,-v)
\]
preserves the rectangle and interchanges the upper forbidden corner
with the lower forbidden corner $b+v<-m$. The two corners are disjoint.
Consequently the fixed-$s$ slice has
\begin{equation}\label{eq:four-slice-polynomial}
 F_n(m)=(n+1-\epsilon)(2m+1)-h(h+1)
\end{equation}
points. The central slice $s=0$ corresponds to $m=n$, and all others
are paired by negation, giving
\begin{equation}\label{eq:four-slices}
 A_4(n)=F_n(n)+2\sum_{m=0}^{n-1}F_n(m).
\end{equation}
In particular,
\[
 F_n(2q)=-q^2+(4n+3)q+n+1,\qquad
 F_n(2q+1)=-q^2+(4n-1)q+3n.
\]
For completeness, the two prefix sums used to evaluate
\eqref{eq:four-slices} are
\begin{align*}
 3\sum_{m=0}^{2q-1}F_n(m)
   &=-2q^3+(12n+6)q^2-q,\\
 3\sum_{m=0}^{2q}F_n(m)
   &=-2q^3+(12n+3)q^2+(12n+8)q+3n+3.
\end{align*}
Each follows by summing the displayed quadratic slice polynomials,
or by induction on $q$. Substitution of $n=2q$ and $n=2q+1$ in
\eqref{eq:four-slices} yields the asserted parity formula.
\end{proof}

\begin{theorem}\label{thm:bounded-zero-sum-gf}
As identities of formal power series,
\begin{align*}
 \sum_{n\ge1}A_3(n)x^n
   &=\frac{x(3+3x+2x^2-x^3+x^4)}{(1-x)^3(1+x+x^2)},\\
 \sum_{n\ge1}A_4(n)x^n
   &=\frac{x(7+10x+2x^2+4x^3-x^4)}{(1-x)^4(1+x)}.
\end{align*}
For every $n\ge6$ in the source's indexing,
\begin{align*}
 A_3(n)&=2A_3(n-1)-A_3(n-2)+A_3(n-3)
                    -2A_3(n-4)+A_3(n-5),\\
 A_4(n)&=3A_4(n-1)-2A_4(n-2)-2A_4(n-3)
                    +3A_4(n-4)-A_4(n-5).
\end{align*}
The same recurrences hold at $n=5$ with the natural extension
$A_3(0)=A_4(0)=1$.
\end{theorem}

\begin{proof}
The formulas of Theorem~\ref{thm:bounded-zero-sum} give
\[
 A_3(n+3)-A_3(n)=8n+16,\qquad
 A_4(n+2)-A_4(n)=11n^2+33n+30.
\]
Taking two further finite differences in the first identity and three
in the second proves the stated recurrences. The corresponding
expanded denominator polynomials are
\begin{align*}
 D_3(x)&=1-2x+x^2-x^3+2x^4-x^5
          =(1-x)^3(1+x+x^2),\\
 D_4(x)&=1-3x+2x^2+2x^3-3x^4+x^5
          =(1-x)^4(1+x).
\end{align*}
Write $S_r(x)=\sum_{n\ge1}A_r(n)x^n$, whose constant coefficient is
zero. The recurrences imply $[x^n](D_rS_r)=0$ for all $n\ge6$.
The initial positive-index terms are
\[
 (A_3(1),\ldots,A_3(5))=(3,9,17,27,41),\qquad
 (A_4(1),\ldots,A_4(5))=(7,31,81,171,309).
\]
They give
\[
 D_3S_3=3x+3x^2+2x^3-x^4+x^5,\qquad
 D_4S_4=7x+10x^2+2x^3+4x^4-x^5.
\]
Both denominators have constant coefficient one, so they are
invertible as formal power series; division yields the claims.
\end{proof}

\paragraph{Formal scope and prior work.}
The Lean development defines each sequence as the cardinality of the
literal finite set of source arrays. It proves the signed-slice
bijections, both triangular corner counts, the disjoint partition,
the polynomial sums, and coefficientwise generating-function
identities. The formal generating functions use coefficient zero
equal to zero and source counts at every positive index. All seven
modules are default build roots under pinned Lean~4.19.0 and mathlib.
The floor sequence underlying the three-entry answer is already
A143978 \cite{a202253-floor}; its elementary algebra is prior
background. The existing open-math-agents A202253 file is an admitted
statement, and Sequencelib supplies a synthetic formula definition;
neither inspected file proves the source array count
\cite{a202253-oma,a202253-sequencelib}. This package treats two entries
as one family, with recurrences and generating functions as
consequences of the count formulas. A202253's separate empirical
interpretation by bounded solutions of $w+2x+3y=0$ is outside the
selected claims. The bounded searches recorded with the source
snapshots do not establish worldwide priority.
